\setupbodyfont[10pt]
\setuppapersize[A4]
\definelinenumbering[ln][step=1,location=left]
\def\Q#1#2{\blank\noindent{\bf QUESTÃO #1}\par #2\par
\noindent A) primeira alternativa da questão\par
\noindent B) segunda alternativa da questão\par
\noindent C) terceira alternativa da questão\par
\noindent D) quarta alternativa da questão\par
\noindent E) quinta alternativa da questão\par}
\starttext
\Q{01}{Enunciado da questão um, com texto suficiente para a extração.}
\Q{02}{Enunciado da questão dois, com texto suficiente para a extração.}
\Q{03}{Enunciado da questão três, com texto suficiente para a extração.}
\Q{04}{Enunciado da questão quatro, com texto suficiente para a extração.}
\Q{05}{Enunciado da questão cinco, com 50\% de chance & custo de R\$ 10.}
\blank
\noindent{\bf Texto para as Questões 06 a 10.}\par
\midaligned{\bf De próprio punho}
\midaligned{\it A escrita e suas tecnologias sofrem interessantes metamorfoses}
\startlinenumbering[ln]
\startlines
Estranhei muito na primeira vez que escutei a expressão “de próprio punho”.
Parecia que eu ia bater em alguém. Não era bem o caso. Foi numa situação
bancária, dessas bem burocráticas, e eu devia escrever algo bem breve,
mas com minhas mãos. Na verdade, o que importava era a autenticidade
da minha caligrafia, que à época ainda era mais fluente e firme.
\stoplines
\stoplinenumbering
\Q{06}{No que diz respeito ao gênero bilhete, a autora dessa crônica}
\Q{07}{O elemento que caracteriza esse texto como crônica é a}
\Q{08}{Nesse texto, a escrita digital é ilustrada pelo(a)}
\Q{09}{A autora conclui que as novas tecnologias de escrita}
\Q{10}{O recurso linguístico usado na síntese da opinião (l. 3-4) foi o(a)}
\page
\startcolumns[n=2]
\Q{11}{Questão onze, antes do texto comum, na coluna da esquerda.}
\blank
\noindent{\bf Texto para os itens 12 e 13}\par
Minha terra tem palmeiras, onde canta o sabiá; as aves que aqui gorjeiam não gorjeiam como lá. Nosso céu tem mais estrelas, nossas várzeas têm mais flores.
\blank
\noindent{\bf QUESTÃO 12}\par Julgue o item: o poema exalta a terra natal. Certo ou errado.\par
\noindent{\bf QUESTÃO 13}\par Julgue o item: o poema usa versos livres. Certo ou errado.\par
\stopcolumns
\stoptext
